\chapter{Bug 修复结构演化分析（静态）}\label{chap:static}

\section{AST 与 libcst 分析方法}
本章目标是回答“修复代码到底在结构上改变了什么”。为避免仅停留在文本 diff 层，本研究采用“\textbf{diff 抽取 $\rightarrow$ 结构特征提取 $\rightarrow$ 模式归纳}”的三步流程，并与仓库已有产物对齐：
\begin{enumerate}
    \item 数据准备：使用 \texttt{static\_analyze/diff\_extract.py} 生成 677 份补丁差异文件（\texttt{static\_analysis\_diffs/*.diff}）；
    \item 结构提取：基于 AST/语义规则识别 \texttt{if/else}、谓词组合、错误路径与返回路径变化，输出到 \texttt{ast\_logic\_analysis\_report.csv}；
    \item 语义归纳：结合 \texttt{bug\_fix\_structure\_feature.md} 与样本 diff，统一抽象为跨案例可迁移的修复模式。
\end{enumerate}

需要强调的是，PostgreSQL 代码中宏与类型别名密集，直接解析 C AST 容易出现语法中断；因此分析中采取了预处理脱敏策略，先保证可解析性，再统计节点差异。该步骤并不改变修复语义，只降低语法噪声对结构统计的干扰。

\subsection*{样本范围与覆盖}
样本覆盖 2022--2025 年，且以中等规模补丁为主，以减少“纯格式化”与“大规模重构”带来的偏差。路径聚合统计显示 677 条样本中，\texttt{storage} 与 \texttt{replication} 各 45 条，\texttt{executor} 31 条，\texttt{catalog} 24 条，\texttt{planner} 22 条，\texttt{wal} 13 条。该分布与前文“高风险模块长期活跃”的观察一致。

\section{典型修复模式：前置条件显式化与错误路径显式化}
通过对 677 条样本进行结构归并，本研究识别出两类最稳定、最具工程解释力的模式。

\subsection*{模式 A：前置条件显式化（Precondition Explicitness）}
该模式的核心是把“默认成立”的假设改写成“代码可检查”的约束，典型表现包括：
\begin{itemize}
    \item 新增 \texttt{NULL/边界/状态} 校验；
    \item 将主逻辑包裹入 guard，非法输入提前返回；
    \item 将隐含前提从注释语义提升为可执行条件。
\end{itemize}

结构上，这类修复会带来 \texttt{if} 节点数增加、return 路径增加、入口条件变强。它的工程价值在于：将“运行后才暴露的问题”前移为“入口可拒绝的问题”，降低错误传播深度。

\subsection*{模式 B：错误路径显式化（Error Path Materialization）}
该模式对应“silent failure 治理”，表现为：
\begin{itemize}
    \item 从隐式失败切换为显式错误分支；
    \item 新增 \texttt{ereport/elog} 或错误码返回；
    \item 正常路径与异常路径分离，便于测试与审计。
\end{itemize}

从控制流视角看，这一变化等价于“把原本不可见的失败状态加入状态机”。在审计实践中，这类改动可显著提升问题可诊断性，也能为后续形式化规约提供清晰的失败语义。

\subsection*{复合模式：Logic + Memory}
仓库静态报告指出，逻辑分支缺陷常与内存/资源问题联动出现。具体而言，条件遗漏会放大到资源生命周期错误（如释放时机错误、越界触发窗口扩大）。因此，修复并非“局部补丁”，而是对条件空间与副作用边界的共同收紧。

\section{控制流与条件复杂度变化}
本节从“控制流图变化”和“谓词复杂度变化”两个维度细化静态证据。

\subsection*{控制流图（CFG）变化趋势}
修复前，部分样本呈单路径直通执行；修复后常出现以下变化：
\begin{enumerate}
    \item 新增入口 guard 分支，非法状态快速退出；
    \item 将异常流程独立为错误路径，避免与正常路径混杂；
    \item return 点位增加，但路径语义更可解释。
\end{enumerate}

虽然 CFG 节点数可能上升，但这种“复杂度上升”在工程上通常是正向复杂度：系统用更明确的路径表达不变量，而非依赖隐式假设。

\subsection*{条件表达式复杂度变化}
在谓词层面，修复常由“单条件”演化为“复合条件”：
\begin{itemize}
    \item 从 \texttt{x > 0} 扩展为 \texttt{x != NULL \&\& x > 0}；
    \item 从状态单点判断扩展为状态 + 资源 + 上下文联合约束；
    \item 引入 \texttt{AND/OR} 组合后，路径判定边界更清晰。
\end{itemize}

这种演化意味着约束空间被细分，错误输入更难穿透到核心逻辑区，是“防御式编程增强”的直接结构证据。

\subsection*{可维护性影响讨论}
修复后的代码行数与分支数可能增加，但维护成本并不一定恶化。相反，当条件语义、失败路径、日志语义三者对齐时，后续维护者可更快定位问题并进行局部验证。换言之，\textbf{可解释性提升}抵消了部分表面复杂度。

\section{静态分析小结}
本章结论可概括为三条：
\begin{enumerate}
    \item PostgreSQL 修复呈现稳定的结构演化方向：\textbf{约束前移、错误显式化、谓词增强}；
    \item 高风险模块（\texttt{storage/replication/WAL}）中此类模式更集中，说明关键路径治理依赖结构性修复；
    \item 静态分析不仅解释“代码改了哪里”，更输出了可被动态复现与形式化验证复用的条件框架。
\end{enumerate}

因此，静态证据在本研究中承担“模式提炼与规约生成”的基础角色，为第四章运行时验证与第五章证明级验证提供可转译输入。